\documentclass[12pt, a4paper]{article}

\usepackage[T2A]{fontenc}
\usepackage[utf8]{inputenc}
\usepackage[russian]{babel}
\usepackage[left=22mm, right=10mm, top=15mm, bottom=25mm]{geometry}
\usepackage[pages=all]{background}
\usepackage{tikz}
\usepackage{microtype}
\usepackage{rotating, graphicx}
\usepackage{mwe} 
\usepackage{wrapfig}
\usepackage{caption}
\usepackage{float}
\usepackage[hidelinks]{hyperref}

\setlength{\parindent}{1.0cm}

% ===== ПРОСТАЯ РУССКАЯ НУМЕРАЦИЯ БЕЗ ПАКЕТОВ =====
\makeatletter
\def\rusenum#1{%
    \expandafter\ifcase\csname c@#1\endcsname
    \or а%
    \or б%
    \or в%
    \or г%
    \or д%
    \or е%
    \or ё%
    \or ж%
    \or з%
    \or и%
    \or й%
    \or к%
    \or л%
    \or м%
    \or н%
    \or о%
    \or п%
    \or р%
    \or с%
    \or т%
    \or у%
    \or ф%
    \or х%
    \or ц%
    \or ч%
    \or ш%
    \or щ%
    \or ъ%
    \or ы%
    \or ь%
    \or э%
    \or ю%
    \or я%
    \fi
}
\renewcommand{\labelenumi}{\rusenum{enumi})}
\renewcommand{\theenumi}{\rusenum{enumi}}
\makeatother




\newcommand{\MyCustomFrame}{%
    \begin{tikzpicture}[remember picture, overlay]
        % 1. ОСНОВНАЯ РАМКА
        \draw [line width=1pt] ([shift={(20mm,-5mm)}]current page.north west) rectangle ([shift={(-5mm,5mm)}]current page.south east);
        
        % 2. НИЖНИЙ ШТАМП (185x15 мм)
% --- ОБНОВЛЕННЫЙ БЛОК НИЖНЕГО ШТАМПА ---
        \begin{scope}[shift={([shift={(-5mm,5mm)}]current page.south east)}]
            \draw [line width=1pt] (0,0) rectangle (-185mm, 15mm);
            
            % ВЕРТИКАЛЬНЫЕ ЛИНИИ (Расширяем ячейки)
            \draw [line width=0.5pt] (-10mm, 0) -- (-10mm, 15mm);   % Линия для "Лист"
            \draw [line width=0.5pt] (-130mm, 0) -- (-130mm, 15mm); % Правая граница блока с изменениями
            \draw [line width=0.5pt] (-140mm, 0) -- (-140mm, 15mm); % Граница Дата (10мм)
            \draw [line width=0.5pt] (-150mm, 0) -- (-150mm, 15mm); % Граница Подп. (10мм)
            \draw [line width=0.5pt] (-165mm, 0) -- (-165mm, 15mm); % Граница № док. (СДЕЛАЛ 15мм!)
            \draw [line width=0.5pt] (-175mm, 0) -- (-175mm, 15mm); % Граница Лист (10мм)
            % Крайняя левая граница (Изм.) остается -185мм (тоже 10мм шириной)
            
            % Горизонтальные линии
            \draw [line width=0.5pt] (0, 7.5mm) -- (-10mm, 7.5mm);
            \draw [line width=0.5pt] (-130mm, 5mm) -- (-185mm, 5mm);
            \draw [line width=0.5pt] (-130mm, 10mm) -- (-185mm, 10mm);
            
            % Текст "Лист" справа
            \node at (-5mm, 11.25mm) {\scriptsize \textsl{Лист}};
            \node at (-5mm, 3.75mm) {\scriptsize \textsl{\thepage}};
            
            % Шифр (сместил чуть-чуть из-за расширения таблицы)
            \node at (-70mm, 7.5mm) {\huge \textsl{ТПЖА.000000.000 ПЗ}};
            \node[anchor=north east] at (0, -2mm) {\scriptsize \textsl {Формат А4}};
            \node[anchor=north east] at (-6, -2mm) {\scriptsize \textsl {Копировал}};
            
            % ПОДПИСИ (Координаты X теперь ровно посередине новых ячеек)
            \node at (-180mm, 2.5mm) {\scriptsize \textsl{Изм.}};
            \node at (-170mm, 2.5mm) {\scriptsize \textsl{Лист}};
            \node at (-157.5mm, 2.5mm) {\scriptsize \textsl{№ докум.}}; % Центр широкой ячейки
            \node at (-145mm, 2.5mm) {\scriptsize \textsl{Подп.}};
            \node at (-135mm, 2.5mm) {\scriptsize \textsl{Дата}};
        \end{scope}
        % 3. БОКОВАЯ ТАБЛИЦА
        \begin{scope}[shift={([shift={(8mm,5mm)}]current page.south west)}]
            \draw [line width=1pt] (0,0) rectangle (12mm, 145mm);
            \draw [line width=0.5pt] (6mm, 0) -- (6mm, 145mm);
            \draw [line width=0.5pt] (0, 30mm) -- (12mm, 30mm);
            \draw [line width=0.5pt] (0, 61mm) -- (12mm, 61mm);
            \draw [line width=0.5pt] (0, 90mm) -- (12mm, 90mm);
            \draw [line width=0.5pt] (0, 120mm) -- (12mm, 120mm);
            
            \node[rotate=90, anchor=center] at (3.5mm, 14.7mm) {\fontsize{11pt}{12pt} \textsl{Инв. № подл.}};
            \node[rotate=90, anchor=center] at (3.5mm, 44.5mm) {\fontsize{11pt}{12pt} \textsl{Подп. и дата}};
            \node[rotate=90, anchor=center] at (3.5mm, 75.5mm) {\fontsize{11pt}{12pt} \textsl{Взам. инв. №}};
            \node[rotate=90, anchor=center] at (3.5mm, 105mm) {\fontsize{11pt}{12pt} \textsl{Инв. № дубл.}};
            \node[rotate=90, anchor=center] at (3.5mm, 132.5mm) {\fontsize{11pt}{12pt} \textsl{Подп. и дата}};
        \end{scope}
    \end{tikzpicture}%
}
\backgroundsetup{
    scale=1,
    color=black,
    opacity=1,
    angle=0,
    contents={\MyCustomFrame} % Просто вызываем нашу команду
}

\newcounter{game_cnt}
\setcounter{game_cnt}{1}

\newcounter{img_cnt}
\setcounter{img_cnt}{1}


\begin{document}
    \pagestyle{empty}
    \sloppy
    \NoBgThispage
\begin{center}
    {\normalsize Реферат}
\end{center}

Матешко В.В. Разработка однопользовательской 2D игры <<Cuppa>>. ТПЖА.090301.671 ПЗ: Курс. проект / ВятГУ, каф. ЭВМ; рук. Сергеева Е.Г. - Киров, 2026. -- ПЗ с 00, рис. 0, прил. 0

\vspace{18pt}

2D-ИГРА, РИТМ-ИГРА, РАЗРАБОТКА ИГРЫ, ВИЗУАЛЬНАЯ НОВЕЛЛА, СЮЖЕТНОЕ ВЕТВЛЕНИЕ, СИСТЕМА ОЦЕНКИ, ПАРСИНГ SM-ФАЙЛОВ, ОБРАБОТКА ДЛИННЫХ НОТ, СИНХРОНИЗАЦИЯ АУДИО И ГЕЙМПЛЕЯ, GODOT, GDSCRIPT.\\

\vspace{18pt}

\textbf{Объект курсового проекта} -- однопользовательские 2D-игры, музыкальный элемент как ключевая игровая механика, интерактивное повествование, переходы между игровыми состояниями.\\

\textbf{Предмет курсового проекта} -- однопользовательская 2D-игра <<Cuppa>>, сочетающая визуальную новеллу и ритм-механику приготовления кофе, а также проектирование и реализация соответствующих алгоритмов.\\

\textbf{Цель курсового проекта} -- разработка и реализация однопользовательской 2D-игры <<Cuppa>> на игровом движке Godot, в которой точность прохождения ритм-сцен влияет на повествование развития истории через диалоги и итоговые концовки.\\

Результатом выполнения курсового проекта является функциональный прототип игры <<Cuppa>>, включающий главное меню, пролог, диалоговые сцены, ритм-модуль с обработкой обычных и длинных нот, систему оценки попаданий, экран концовки и переходы между всеми ключевыми состояниями.\\

\newpage


\textbf{\normalsize Содержание}
\vspace{24pt}

\noindent\makebox[\linewidth][l]{\hyperref[sec:intro]{\textbf{Введение}}\hfill \textbf{\pageref{sec:intro}}}\\[1.2em]
\noindent\hyperref[sec:analysis]{\textbf{1\hspace{0.70em}Анализ предметной области}} \dotfill \textbf{\pageref{sec:analysis}}\\
\noindent\hspace*{5mm}\hyperref[sec:analysis-overview]{1.1\hspace{0.70em}Обзор предметной области} \dotfill \pageref{sec:analysis-overview}\\
\noindent\hspace*{5mm}\hyperref[sec:analysis-analogs]{1.2\hspace{0.70em}Анализ аналогов} \dotfill \pageref{sec:analysis-analogs}\\
\noindent\hspace*{10mm}\hyperref[sec:va11]{1.2.1\hspace{0.70em}Анализ игры <<VA-11 Hall-A>>} \dotfill \pageref{sec:va11}\\
\noindent\hspace*{10mm}\hyperref[sec:coffee-talk]{1.2.2\hspace{0.70em}Анализ игры <<Coffee Talk>>} \dotfill \pageref{sec:coffee-talk}\\
\noindent\hspace*{10mm}\hyperref[sec:yunyun]{1.2.3\hspace{0.70em}Анализ игры <<Yunyun Syndrome!? Rhythm Psychosis Demo>>} \dotfill \pageref{sec:yunyun}\\
\noindent\hspace*{10mm}\hyperref[sec:glisynth]{1.2.4\hspace{0.70em}Анализ игры <<Glisynth>>} \dotfill \pageref{sec:glisynth}\\
\noindent\hspace*{5mm}\hyperref[sec:req]{1.3\hspace{0.70em}Требования к разработке} \dotfill \pageref{sec:req}\\
\noindent\hspace*{10mm}\hyperref[sec:req-func]{1.3.1\hspace{0.70em}Функциональные требования} \dotfill \pageref{sec:req-func}\\
\noindent\hspace*{10mm}\hyperref[sec:req-nonfunc]{1.3.2\hspace{0.70em}Нефункциональные требования} \dotfill \pageref{sec:req-nonfunc}\\
\noindent\hspace*{5mm}\hyperref[sec:task]{1.4\hspace{0.70em}Постановка задач} \dotfill \pageref{sec:task}\\

\newpage

    \phantomsection\label{sec:intro}
    {\textbf{\normalsize Введение}}
    \vspace{24pt}

На протяжение всего времени развития человечества у людей существовал интерес и жизненная необходимость в развлечениях. Одной из таких ветвей развлечений на  сегодняшний день является индустрия игр. Развитие технологий позволило создавать видеоигры, доступ к которым имеют большиннство людей на Земле.\\

Видеоигры активно пользуются спросом. Людям интересны как и AAA-проекты, так и инди-игры, разрабатываемые от разработчиков-одиночек до команд энтузиастов.\\

Каждый в этом отрасли старается предложить игрокам новые форматы взаимодействия, оригинальные истории, дизайны и уникальные механики, которые впоследствии могут даже привести к открытию нового жанра игр.\\

И все же существуют некоторые жанры, которые оказались практически позабыты разработчиками. Одним из подобных жанров являются ритм-игры. Их суть заключается в попадании по <<нотам>> -- движущимся объектам -- в ритм музыки, что создает интересное взаимодействия игры с пользователем.\\

Таким образом, у меня возникла идея создать игру в жанре ритм-игры, тем самым внеся свой вклад в развитие данного жанра. \\

Однако я преследую цель создать уютную игру, в которую может поиграть человек в плохом настроение, после работы или просто желающий получить душевное умиротворение. Чтобы создать подобную атмосферу я пришла к идеи разработать визуальную-новеллу про вампирскую кофейню, где главный герой -- бариста. В эту визуальную новеллу встроен элемент ритм-игры, отвечающий за приготовление кофе.\\

\textbf{Цель проекта: } разработать оригинальную 2D-игру, в которой присутствует простой сюжет и взаимодействие с окружением. В ходе игры пользователь будет общаться с вампирами и готовить кофе клиентам, играя в ритм-игру, тем самым зарабатывая очки, от которых зависит концовка игра -- будет ли существовать вампирская кофейня или нет. \\

\textbf{Задачи проекта: }

\begin{enumerate}
    \item Проанализировать существующие проекты и выделить механики, применимые в рамках выбранного жанра;
    \item Сформировать технические и геймдизайнерские требования к проекту;
    \item Спроектировать архитектуру игры, определив структуру классов и состояний;
    \item Разработать базовые модули и реализовать их с использованием игрового движка Godot;
    \item Провести тестирование, проанализировать результаты и определить потенциал дальнейшего развития.
\end{enumerate}

\textbf{Объект исследования: } однопользовательские 2D-игры, музыкальная индустрия, акустические вычисления, игровые механики ритм-игр, игровые системы управления действиями персонажа, состояния объектов, взаимодействие с окружающей средой и пользовательский опыт.\\

\textbf{Предмет исследования: } однопользовательская визуальная новелла <<Cuppa>> с элементами ритм-игры, алгоритмы акустического анализа, динамическая синхронизация игровых событий с музыкой, проектирование и реализация уровней ритм-игры, а также разработка и реализация игрового сюжета, мира и диалогов между персонажами. \\

\textbf{Актуальность проекта} обусловлена сочетанием сразу нескольких факторов:

\begin{enumerate}
    \item Рост интереса к инди-проектам;
    \item Рост интерес к ритм-играм, в том числе с появлением в этом году новых проектов от известных разработчиков; 
    \item Повышенный спрос на медитативные и уютные игры;
    \item Повышенный интерес и любовь в кофе, а также романтизация работы баристы и вампиров;
    \item Интерес к нестандартному совмещению жанров, так как это создает ощущение чего-то нового.
\end{enumerate}

\newpage

    \phantomsection\label{sec:analysis}
    {\textbf{\normalsize 1 Анализ предметной области}}
    \vspace{24pt}

\textbf{Ритм-игры} -- это жанр видеоигр, в основе которого лежит точное следование ритму музыкального сопровождения игроком посредством своесвоевременного выполнения действий, что влияет на набор очков, качество прохождения и итоговый результат.\\

\textbf{Визуальная новелла} -- это жанр видеоигр, игровой процесс которого заключается в чтение текста, выбора реплик или действий, влияющих на развитие истории и её концовку.\\

    \phantomsection\label{sec:analysis-overview}
    \vspace{24pt}
    {\textbf{\normalsize 1.1 Обзор предметной области}}
    \vspace{24pt}

Ритм-игры представляют собой жанр видеоигр, в основе которого лежит взаимодействие игрока с музыкальным сопровождением через точное выполнение действий в заданный ритм. Игровой процесс строится на повторяющихся механиках и синхронизации с музыкой, что формирует у игрока ощущение вовлечённости и постепенного прогресса.\\

Несмотря на кажущуюся простоту, ритм-игры требуют концентрации, реакции и чувства ритма. При этом сам процесс остаётся достаточно понятным и интуитивным, благодаря чему такие игры доступны широкой аудитории.\\

Однако в большинстве ритм-игр основной акцент делается именно на механике, в то время как сюжет и атмосфера либо сильно упрощены, либо вовсе отсутствуют. В результате игровой опыт часто ограничивается только выполнением действий под музыку.\\

В рамках проекта «Cuppa» предлагается объединение ритм-игры и визуальной новеллы. Такой подход позволяет не только сохранить основную механику жанра, но и дополнить её повествованием, персонажами и взаимодействием с игровым миром.\\

Визуальные новеллы, в свою очередь, ориентированы на сюжет, диалоги и развитие истории через выборы игрока. Они создают эмоциональную привязанность к персонажам и позволяют глубже погрузиться в атмосферу игры.\\

Сочетание этих двух жанров даёт возможность создать более разнообразный игровой процесс, где механика ритм-игры встроена в контекст происходящего. В данном проекте это реализуется через концепцию вампирской кофейни, где игрок выступает в роли бариста, взаимодействует с посетителями и готовит напитки, используя элементы ритм-игры.\\

Таким образом, проект направлен на создание не только игрового, но и эмоционального опыта, сочетающего в себе интерактивность, атмосферу и чувство уюта.\\

\newpage

    \phantomsection\label{sec:analysis-analogs}
    \vspace{24pt}
    {\textbf{\normalsize 1.2 Анализ аналогов}}
    \vspace{24pt}

Для выбора архитектурных решений, проектирования игровых механик и определения целевой аудитории важно обратиться к анализу существующих аналогов, что поможет лучше понять, какие подходы уже используются разработчиками и какие из них можно адаптировать в рамках собственного проекта.\\

В рамках данного исследования рассматриваются игры, обладающие преимущественно следующими признаками:

\begin{enumerate}
    \item Наличие главного героя, выступающего в роли обсулживающего персонала (например, бариста, официант или бармен);
    \item Визуальное исполнение в 2D стиле;
    \item Принадлежность к жанру визуальной новеллы, ритм-игры или их сочетанию;
    \item Акцент на взаимодействия с персонажами и элементами окружения.
\end{enumerate}

    \phantomsection\label{sec:va11}
    \vspace{24pt}
    {\textbf{\normalsize 1.2.1 Анализ игры <<VA-11 Hall-A>>}}
    \vspace{24pt}

Игра представляет собой визуальную новеллу, в которой игрок выступает в роли бармена в киберпанк-баре. Основной игровой процесс строится вокруг общения с посетителями и приготовления напитков, от которых зависит развитие диалогов и сюжетных линий.\\

Особенностью проекта является то, что механика приготовления напрямую связана с повествованием: выбор ингредиентов влияет на состояние персонажей и их поведение. Визуально игра выполнена в 2D-пиксель-арте с акцентом на атмосферу и стиль.\\

\textbf{Преимущества:}

\begin{itemize}
    \item Интересное сочетание геймплея и повествования;
    \item Сильная атмосфера и проработанные персонажи;
    \item Влияние действий игрока на развитие диалогов.
\end{itemize}

\textbf{Недостатки:}

\begin{itemize}
    \item Ограниченность игровых механик;
    \item Невысокая реиграбельность при повторном прохождении.
\end{itemize}

    \begin{figure}[H]
        \centering
        \includegraphics[scale=0.4]{game\arabic{game_cnt}.png}\\
        \caption*{Рисунок \arabic{img_cnt} -- Пример игрового процесса <<VA-11 Hall-A>>}
        \stepcounter{img_cnt}
        \stepcounter{game_cnt}
    \end{figure}

    \phantomsection\label{sec:coffee-talk}
    \vspace{24pt}
    {\textbf{\normalsize 1.2.2 Анализ игры <<Coffee Talk>>}}
    \vspace{24pt}

Игра представляет собой визуальную новеллу, в которой игрок выступает в роли бариста в кофейне. Основной игровой процесс включает приготовление напитков и общение с посетителями, каждый из которых имеет свою историю и характер.\\

Проект делает акцент на создании уютной и спокойной атмосферы. Механика приготовления напитков упрощена, но органично встроена в игровой процесс и поддерживает общее настроение игры.\\

\textbf{Преимущества:}

\begin{itemize}
    \item Уютная и медитативная атмосфера;
    \item Простая и понятная игровая механика;
    \item Акцент на диалогах и взаимодействии с персонажами.
\end{itemize}

\textbf{Недостатки:}

\begin{itemize}
    \item Ограниченная глубина геймплея;
    \item Небольшое влияние игрока на развитие сюжета.
\end{itemize}

     \begin{figure}[H]
         \centering
         \includegraphics[scale=0.35]{game\arabic{game_cnt}.png}\\
         \caption*{Рисунок \arabic{img_cnt} -- Пример игрового процесса <<Coffee Talk>>}
         \stepcounter{img_cnt}
         \stepcounter{game_cnt}
     \end{figure}

    \phantomsection\label{sec:yunyun}
    \vspace{24pt}
    {\textbf{\normalsize 1.2.3 Анализ игры <<Yunyun Syndrome!? Rhythm Psychosis Demo>>}}
    \vspace{24pt}

Игра представляет собой ритм-игру с ярким визуальным стилем и необычной подачей. Игровой процесс интерпретируется как внутреннее состояние главной героини, где музыкальные ритмы отражают психоз и попытки справиться с реальностью.\\ 

\textbf{Преимущества:}

\begin{itemize}
    \item Выразительный визуальный стиль;
    \item Динамичный и вовлекающий геймплей;
    \item Аккуратная синхронизация визуала и музыки;
    \item Простая и понятная игровая механика.
\end{itemize}

\textbf{Недостатки:}

\begin{itemize}
    \item Высокая нагрузка на восприятие игрока;
    \item Может быть сложной для неподготовленной аудитории.
\end{itemize}

     \begin{figure}[H]
         \centering
         \includegraphics[scale=0.35]{game\arabic{game_cnt}.png}\\
         \caption*{Рисунок \arabic{img_cnt} -- Пример игрового процесса <<Yunyun Syndrome!? Rhythm Psychosis Demo>>}
         \stepcounter{img_cnt}
         \stepcounter{game_cnt}
     \end{figure}


    \phantomsection\label{sec:glisynth}
    \vspace{24pt}
    {\textbf{\normalsize 1.2.4 Анализ игры <<Glisynth>>}}
    \vspace{24pt}

Игра представляет собой экспериментальную ритм-игру, в которой основной акцент сделан на взаимодействии игрока с аудиовизуальной средой. Игровой процесс строится на синхронизации действий с музыкой и визуальными эффектами.\\

Проект отличается минималистичным подходом к дизайну и фокусом на ощущениях от процесса, а не на классической системе прогрессии.\\

\textbf{Преимущества:}

\begin{itemize}
    \item Оригинальный подход к ритм-механикам;
    \item Сильная связь между звуком и визуалом;
    \item Атмосферность и экспериментальность;
\end{itemize}

\textbf{Недостатки:}

\begin{itemize}
    \item Отсутствие выраженного сюжета;
    \item Ограниченность игровых целей и мотивации.
\end{itemize}

     \begin{figure}[H]
         \centering
         \includegraphics[scale=0.25]{game\arabic{game_cnt}.png}\\
         \caption*{Рисунок \arabic{img_cnt} -- Пример игрового процесса <<Glisynth>>}
         \stepcounter{img_cnt}
         \stepcounter{game_cnt}
     \end{figure}

    \phantomsection\label{sec:req}
    \vspace{24pt}
    {\textbf{\normalsize 1.3 Требования к разработке}}
    \vspace{24pt}

Требования к разрабатываемому программному продукту можно условно разделить на функциональные и нефункциональные. Первые описывают основные возможности системы и игровой процесс, тогда как вторые определяют качество работы, визуальную составляющую и технические ограничения проекта.\\

    \phantomsection\label{sec:req-func}
    {\textbf{\normalsize 1.3.1 Функциональные требования}}

\begin{enumerate}
    \item \textbf{Система диалогов и повествования:} отображение текстовой информации, а также возможность выбора реплик, влияющих на развитие сюжета и взаимодействие с персонажами;
    \item \textbf{Ритм-механика приготовления напитков:} реализация игрового процесса, основанного на попадании в ритм музыкального сопровождения;
    \item \textbf{Обработка игровых событий:} генерация <<нот>> на основе заданных данных;
    \item \textbf{Система оценки и прогрессии:} определение точности действий игрока с последующим начислением очков, влияющих на качество приготовления напитков и концовку игры;
    \item \textbf{Игровое окружение кофейни:} реализация взаимодействия с посетителями по схеме «заказ -- приготовление -- отдача», связывающей элементы визуальной новеллы и ритм-игры;
    \item \textbf{Система состояний:} управление переходами между различными режимами игры (меню, диалоги, ритм-игра) без потери текущего прогресса;
    \item \textbf{Пользовательский опыт:} звуковая и визуальная реакция системы на действия игрока;
    \item \textbf{Пользовательский интерфейс:} интуитивно понятный интерфейс, отображающий диалоги, элементы ритм-игры, количество очков и состояние текущего процесса.
\end{enumerate} \

    \phantomsection\label{sec:req-nonfunc}
    {\textbf{\normalsize 1.3.2 Нефункциональные требования}}

\begin{enumerate}
    \item \textbf{Визуальная стилистика:} 2D в стиле пиксель-арт с «ночной» атмосферой и яркими акцентами, создающими ощущение уюта;
    \item \textbf{Музыкальное сопровождение:} спокойная музыка в стиле lo-fi и джаза  с чётко выраженным ритмом;
    \item \textbf{Синхронизация аудио и визуала:} минимальная задержка между звуком и отображением элементов для комфортного игрового процесса;
    \item \textbf{Производительность:} стабильная работа игры на персональном компьютере без заметных просадок;
    \item \textbf{Надёжность:} устойчивость системы к действиям пользователя, способным привести к сбоям;
    \item \textbf{Читаемость текста:} использование элегантных шрифтов, подходящих стилю вампиров, при этом сохраняя удобство восприятия.
\end{enumerate}

    \phantomsection\label{sec:task}
{\textbf{\normalsize 1.4 Постановка задач}}

\end{document}
